\documentclass[11pt,a4paper]{article}
\usepackage[margin=25mm,headheight=15pt]{geometry}
\usepackage[T1]{fontenc}
\usepackage{lmodern}
\usepackage{microtype}
\usepackage{amsmath,amssymb,amsthm,mathtools}
\usepackage{booktabs,longtable,array}
\usepackage[dvipsnames]{xcolor}
\usepackage{enumitem}
\usepackage{fancyhdr}
\usepackage{hyperref}
\usepackage{xurl}
\hypersetup{colorlinks=true,linkcolor=MidnightBlue,urlcolor=MidnightBlue,citecolor=MidnightBlue,
 pdftitle={Coefficient height one and an eleven-integer degree-464 equation for Con(ZFC)},
 pdfauthor={OpenAI-assisted research prepared for Vladimir Reshetnikov}}
\setlist{topsep=4pt,itemsep=3pt,parsep=0pt}
\setlength{\parskip}{4pt}
\setlength{\parindent}{0pt}
\setlength{\emergencystretch}{2em}
\pagestyle{fancy}
\fancyhf{}
\fancyhead[L]{\small\textsc{Diophantine consistency equations}}
\fancyhead[R]{\small 24 October 2026}
\fancyfoot[C]{\thepage}
\newtheorem{theorem}{Theorem}[section]
\newtheorem{lemma}[theorem]{Lemma}
\newtheorem{proposition}[theorem]{Proposition}
\newtheorem{corollary}[theorem]{Corollary}
\theoremstyle{definition}
\newtheorem{definition}[theorem]{Definition}
\theoremstyle{remark}
\newtheorem{remark}[theorem]{Remark}
\newcommand{\ZFC}{\mathsf{ZFC}}
\newcommand{\ZF}{\mathsf{ZF}}
\newcommand{\Con}{\operatorname{Con}}
\newcommand{\HeightOneVariables}{7431}
\newcommand{\file}[1]{\nolinkurl{#1}}
\newcommand{\Height}{\operatorname{ht}}
\DeclareMathOperator{\sgn}{sgn}
\title{\vspace{-10mm}\textbf{Coefficient height one and an eleven-integer\\[3pt]
degree-464 equation for $\Con(\ZFC)$}}
\author{Research prepared for Vladimir Reshetnikov and the ProveIt project\\[3pt]
\small OpenAI-assisted mathematical development and independent agent review}
\date{24 October 2026}

\begin{document}
\maketitle
\thispagestyle{empty}
\begin{abstract}
We give two explicit ordinary polynomial equations over unrestricted
integers whose lack of solutions is provably equivalent in ZFC to
syntactic consistency of ZFC, using the explicitly retained calibrations
of the ProveIt repository. The first is a fully expanded quartic with
7,431 variables and 34,086 nonzero coefficients, every one equal to $1$
or $-1$. Its constant coefficient is one, and its individual variable
degrees are at most two. A general block-and-cycle construction gives
an exact bijection between its zero set and that of a supplied direct
C222 quadratic system. This route does not simulate K215 and does not
instantiate Jones's universal polynomial.

The second equation has eleven variables, exact total degree 464, and
a division-free circuit of 122 live binary arithmetic operations. It
combines selective coprime radical fusion with powered divisibility,
fixed-pair Pell interpolation, a reduced guard, and a shared auxiliary
center. A direct companion has degree 488 in 121 operations. The
eleven-variable construction retains the complete fixed native
coefficient tuple. Both reductions have ordinary mathematical proofs;
their finite certificates, source bytes, exact polynomial data, and
reproducible checkers are supplied. No global simplicity minimum,
worldwide priority, proof of consistency, or new end-to-end formal
proof-assistant certificate is asserted.
\end{abstract}

\vspace{3mm}
\begin{center}
\begin{tabular}{lrrr}
\toprule
New equation & Integer variables & Exact degree & Principal size objective\\
\midrule
Height-one quartic $\mathcal H$ & 7,431 & 4 & $\Height(\mathcal H)=1$\\
Centered norm $\mathcal E_{464}$ & 11 & 464 & 122 arithmetic operations\\
Direct norm $\mathcal E_{488}$ & 11 & 488 & 121 arithmetic operations\\
\bottomrule
\end{tabular}
\end{center}

\clearpage
\tableofcontents
\clearpage

\section{Objective, comparison, and logical contract}

The requested target is a fixed polynomial $P\in\mathbb Z[x_1,\ldots,x_n]$
such that
\begin{equation}
 \ZFC\vdash
 \left(\forall x_1,\ldots,x_n\in\mathbb Z\;P(x_1,\ldots,x_n)\ne0\right)
 \longleftrightarrow\Con(\ZFC).
 \label{eq:target}
\end{equation}
Here $\Con(\ZFC)$ means the absence of a finite syntactic proof of
contradiction in the fixed standard proof system. Every variable in this
report is unrestricted unless explicitly used in an intermediate lemma.
Positive, nonnegative, and unrestricted variable counts are not
interchanged.

The classical Diophantine representation theorem establishes the
availability of such equations in principle. Recent work on universal
pairs and reductions over the integers refines the number of unknowns
and degree \cite{bayer,sun}. Those general results do not select the
smallest description for this particular consistency statement. A
source-specific construction must also fix its coefficient tuple or
literal residual list and establish its refutation calibration.

There is no single agreed meaning of ``simplest.'' We measure the vector
\[
 (\text{domain},\ \text{variables},\ \text{degree},\ \text{coefficient height},
 \ \text{support},\ \text{circuit operations},\ \text{coefficient storage}).
\]
The coefficient height used throughout is
$\Height(P)=\max_\alpha|c_\alpha|$ for the \emph{collected} polynomial
$P=\sum_\alpha c_\alpha x^\alpha$. It is not the weighted sum sometimes
also called a height in experimental Diophantine searches. A circuit
operation is a live binary addition, subtraction, or multiplication,
including squares and multiplication by fixed constants; leaves and
reuse are free. This count does not measure bit complexity.

\subsection{Frozen sources}

The comparison snapshot is the exact ProveIt commit
\begin{center}
\small\file{45b399ec51dfbecb891dd51d8a9e4b6b7c361ac6}.
\end{center}
The audit included the active research summaries, the delivered forms
through package 110, and all eight then-present incoming archives.
The incoming archives matter: their results already improve on several
numbers in the collection README. In particular, the direct source
chosen here is the 519-variable equation in the powered-saturation
archive, while the eleven-variable comparison comes from the coprime
projection archive. Their immutable paths and SHA-256 hashes are
retained with the artifacts.

\begin{center}
\begin{tabular}{lrrl}
\toprule
Compared construction & Variables & Degree & Additional metric\\
\midrule
Incoming height-two quartic \cite{heighttwo} & 3,564 & 4 & Height 2; 14,518 terms\\
Incoming direct source \cite{round14} & 519 & 4 & Height 411,060; 6,253 terms\\
New height-one quartic & 7,431 & 4 & Height 1; 34,086 terms\\
\midrule
Incoming centered coprime norm \cite{coprime} & 11 & 688 & 126 operations\\
Incoming direct coprime norm \cite{coprime} & 11 & 728 & 122 operations\\
New centered norm & 11 & 464 & 122 operations\\
New direct norm & 11 & 488 & 121 operations\\
\bottomrule
\end{tabular}
\end{center}
All seven rows use unrestricted integer coordinates. The height-one
construction improves the coefficient objective while increasing
dimension and support. The centered norm lowers degree by
$224/688\approx32.56\%$ and lowers operation count from 126 to 122
relative to its matched centered source. Ten-integer equations and a
native 37-integer quartic already occur in the inspected repository.
Neither new equation is advertised as the smallest by variable count.

\subsection{Three kinds of evidence}

The report distinguishes ordinary unbounded proofs, exact finite
artifact checks, and formal proof-assistant endpoints. The new
mathematics below is proved by ordinary arguments formalizable in ZFC.
The scripts verify literal polynomial identities, coefficients,
dependency graphs, degree attainments, and source bindings. These checks
do not decide whether any equation has an integer zero. No new complete
Lean, Rocq, or object-language ZFC proof has been built in this round.
Inherited formal endpoints are named as part of the provenance rather
than described as newly checked.

\section{A quartic with coefficient height one}\label{sec:heightone}

The principal new mechanism is a coefficient reduction that applies to any
finite quadratic system over the integers. It uses nonnegative quadratic
forms to enforce agreement between separate occurrences of a variable.
The separation is essential: algebraically identifying these occurrences
before collecting coefficients would destroy the height bound.

\subsection{The block identity}

Let $\sigma_1,\ldots,\sigma_k\in\{-1,1\}$. For two lists of real numbers
$m=(m_1,\ldots,m_k)$ and $w=(w_1,\ldots,w_k)$, define
\begin{align}
\mathcal B_\sigma(m,w)
  ={}&\sum_{i=1}^k(m_i^2+w_i^2-m_iw_i)\\[-2mm]
   &+\sum_{1\le i<j\le k}\sigma_i\sigma_j(m_im_j+w_iw_j).
   \label{eq:block}
\end{align}
The exact identity
\begin{equation}
 2\mathcal B_\sigma(m,w)=
 \left(\sum_i\sigma_i m_i\right)^2+
 \left(\sum_i\sigma_i w_i\right)^2+
 \sum_i(m_i-w_i)^2
 \label{eq:block-sos}
\end{equation}
proves both nonnegativity and the equivalence
\begin{equation}
 \mathcal B_\sigma(m,w)=0
 \quad\Longleftrightarrow\quad
 w_i=m_i\ (1\le i\le k),\qquad \sum_i\sigma_i m_i=0.
 \label{eq:block-zero}
\end{equation}
These statements hold over any ordered field, and consequently over
$\mathbb Z$ as well as $\mathbb R$. The factor $1/2$ in the proof is harmless:
the definition \eqref{eq:block} is already an integer polynomial.

Ordinary squaring of a residual produces coefficients $2$ in its mixed
terms. Formula \eqref{eq:block} distributes each diagonal contribution
between two matching constraints. After division by two, every displayed
coefficient is $1$ or $-1$. The remaining issue is to prevent distinct
displayed terms from becoming the same monomial.

\subsection{Separating occurrences and preserving degree four}

First write each source row as a list of signed monomial occurrences,
\begin{equation}
 R_j(x)=\sum_{i=1}^{k_j}\sigma_{ji}M_{ji}(x),
 \qquad \deg M_{ji}\le 2.
 \label{eq:unit-rows}
\end{equation}
A coefficient $c$ can be written as $|c|$ occurrences with sign
$\operatorname{sgn}(c)$. Section~\ref{sec:binary-normalization} gives a
more economical preliminary normalization for large coefficients.
Occurrences remain distinct even when their monomials agree in the
source row.

For each occurrence introduce two entirely fresh variables $a_{ji},b_{ji}$
and set $m_{ji}=a_{ji}b_{ji}$. Label the two factors as follows:
\begin{center}
\begin{tabular}{lll}
\toprule
Source monomial & Label of $a_{ji}$ & Label of $b_{ji}$\\
\midrule
$x_p x_q$ & $x_p$ & $x_q$\\
$x_p$ & $x_p$ & $1$\\
$1$ & $1$ & $1$\\
\bottomrule
\end{tabular}
\end{center}
Even for $x_p^2$, the two aliases are different variables with the same
label. No alias is used in two occurrences. Introduce one further fresh
variable $w_{ji}$ for each occurrence, and use the block
$\mathcal B_{\sigma_j}(m_j,w_j)$.

All $m_{ji}$ have degree exactly two, including padded linear and constant
occurrences. Thus the $m_{ji}m_{j\ell}$ terms have degree four,
the $m_{ji}w_{ji}$ terms have degree three, and the $w_{ji}w_{j\ell}$
terms have degree two. This deliberate separation of degrees is what
allows the agreement constraints below to be added without coefficient
collisions.

\subsection{Agreement cycles and the single unit anchor}

For a list $z_1,\ldots,z_d$ with $d\ge3$, put $z_{d+1}=z_1$ and define
\begin{equation}
 \mathcal C(z_1,\ldots,z_d)
   =\sum_{i=1}^d z_i^2-\sum_{i=1}^d z_i z_{i+1}
   =\frac12\sum_{i=1}^d(z_i-z_{i+1})^2.
 \label{eq:cycle}
\end{equation}
This polynomial is nonnegative and vanishes exactly when all its entries
agree. Its coefficients are $1$ and $-1$. The restriction $d\ge3$ is
necessary: at $d=2$ the two edge monomials coincide and give coefficient
$-2$.

Group all aliases carrying the same old-variable label $x_p$.
For a group of size at least three, add its cycle \eqref{eq:cycle}.
For a group of size two, add one fresh dummy alias with the same label
and use the resulting triangle. A singleton needs no agreement
constraint. An unused source variable can be omitted when only existence
of solutions matters; all source variables of the instantiated equation
occur.

List all aliases with label $1$ as $u_1,\ldots,u_d$. Add fresh unit
aliases if necessary so that $d\ge2$, and put
\begin{align}
 \mathcal U(u_1,\ldots,u_d)
   & =1+\sum_{i=1}^d u_i^2-u_1-u_d
       -\sum_{i=1}^{d-1}u_i u_{i+1}\label{eq:unit-anchor}\\
   & =\frac12\left[(1-u_1)^2+
       \sum_{i=1}^{d-1}(u_i-u_{i+1})^2+(u_d-1)^2\right].
       \nonumber
\end{align}
It vanishes exactly when every $u_i=1$. Only this block contributes a
constant term, namely $1$, or linear terms, namely $-u_1-u_d$.

\begin{definition}[The height-one polynomial]\label{def:heightone}
For the fixed occurrence lists \eqref{eq:unit-rows}, define
\begin{equation}
 \mathcal H=
 \sum_j\mathcal B_{\sigma_j}(m_j,w_j)
   +\sum_{\text{old-variable groups}}\mathcal C
   +\mathcal U.
 \label{eq:heightone}
\end{equation}
The second sum includes the prescribed triangles for groups of size two
and omits singletons. All sums are fixed and finite. All coordinates
range independently over $\mathbb Z$.
\end{definition}

\begin{theorem}[Coefficient and zero-set theorem]\label{thm:heightone}
For a finite quadratic integer system with at least one nonzero monomial
occurrence, the construction above
produces an ordinary polynomial $\mathcal H$ with the following properties:
\begin{enumerate}
\item Every nonzero coefficient of its fully collected expansion is
      $1$ or $-1$; its constant coefficient is $1$.
\item Its total degree is exactly four, and its degree in each individual
      variable is at most two.
\item It is nonnegative on every real tuple.
\item It has an integer zero if and only if all source rows have a
      simultaneous integer zero.
\item If every source variable occurs, projection to one representative
      of each old-variable alias group is a bijection on integer zero
      sets. Its inverse is explicitly polynomial.
\end{enumerate}
\end{theorem}

\begin{proof}
Nonnegativity follows from \eqref{eq:block-sos}, \eqref{eq:cycle}, and
\eqref{eq:unit-anchor}. Suppose that $\mathcal H=0$. Every block is
nonnegative, so every block is zero. The cycles make each old-variable
alias group constant; choose its common value as $x_p$. The anchor makes
all unit aliases equal to one. Consequently each $m_{ji}$ becomes its
source monomial $M_{ji}(x)$. Formula \eqref{eq:block-zero} then gives
$R_j(x)=0$ for every source row.

Conversely, given a simultaneous source zero $x$, set every old-variable
alias to its label's value, every unit alias to one, every agreement
dummy to its group's value, and $w_{ji}=M_{ji}(x)$. Every block vanishes.
All added values are uniquely forced on a zero, which proves the
bijection and its explicit inverse. Before coefficient normalization,
this inverse has degree at most two in the source variables.

For the coefficient claim, partition the collected monomials by degree
and by the variable families they contain. At degree four, the terms
$m_{ji}^2$ use exactly one private alias pair, each with exponent two;
the terms $m_{ji}m_{j\ell}$ use two distinct private pairs. Each pair
belongs to precisely one occurrence, so neither kind can collide with
another term. Across different rows all alias supports are disjoint.

At degree three, each $-m_{ji}w_{ji}$ contains its own private $w$ variable
and its private alias pair. These monomials are all different.
At degree two, the terms containing $w$ variables come only from the row
blocks. A square uses one $w$ index; a mixed term uses a pair from the
same row. They are distinct, and no $w$ is an alias. The remaining
quadratic terms are precisely the agreement cycles and unit anchor.
Their alias groups are disjoint. Every cycle has at least three distinct
vertices, so each undirected edge occurs once. The anchor path likewise
has no repeated edge. Its two linear terms are distinct because $d\ge2$.
The constant term occurs once. Thus every displayed term contributes to
a distinct monomial, proving height exactly one.

Any occurrence contributes $a_{ji}^2b_{ji}^2$ with coefficient one, so
degree four is attained. The same enumeration shows that no individual
exponent exceeds two.
\end{proof}

The polynomial is a rational sum of squares, through the identity for
$2\mathcal H$. An integer sum-of-squares representation of $\mathcal H$
itself is neither used nor asserted. The proof relies on the order of
$\mathbb Z$ or $\mathbb R$; reducing the single equation modulo a prime
does not preserve its zero-set interpretation.

\subsection{A small example: one addition equation}

For $x+y-z=0$, take three unit aliases $u_1,u_2,u_3$, put
$m_1=xu_1$, $m_2=yu_2$, $m_3=zu_3$, and introduce $w_1,w_2,w_3$.
Then
\begin{align*}
\mathcal H_{+}={}&\sum_{i=1}^3(m_i^2+w_i^2-m_iw_i)
 +m_1m_2-m_1m_3-m_2m_3\\
 &+w_1w_2-w_1w_3-w_2w_3
 +1+u_1^2+u_2^2+u_3^2-u_1-u_3-u_1u_2-u_2u_3.
\end{align*}
Its collected coefficients are all $\pm1$, its degree is four, and
$\mathcal H_+=0$ holds exactly when
\[
 x+y=z,\qquad (u_1,u_2,u_3)=(1,1,1),\qquad(w_1,w_2,w_3)=(x,y,z).
\]
This example shows why addition is compatible with height one. No
positivity restriction on $x,y,z$, or on any auxiliary variable, is
needed.

\subsection{Normalizing large coefficients}\label{sec:binary-normalization}

Expanding a coefficient of magnitude $c$ into $c$ occurrences is correct
but can be wasteful. The delivered compiler first constructs each needed
positive integer $c$ from a shared unit and smaller constants. For
$c=2d+\epsilon$ with $\epsilon\in\{0,1\}$, introduce a constant wire
$C_c$ and the row
\begin{equation}
 C_c-C_d-C_d-\epsilon=0.
 \label{eq:binary-constant}
\end{equation}
The repeated $C_d$ is retained as two signed occurrences. Constant wires
are ordered by their defining dependencies. Here $C_1$ denotes the
literal constant $1$ and requires no additional coordinate. The resulting
finite triangular system forces each
$C_c$ to equal its stated integer and has a unique extension of every
source tuple. Sharing the binary ancestors avoids repeating the same
constant construction.

For a coefficient multiplying a quadratic monomial $x_p x_q$, introduce
or reuse a scaled wire $S_{c,p}$ with
\begin{equation}
 S_{c,p}-C_c x_p=0,
 \label{eq:scaled-wire}
\end{equation}
and replace $c x_p x_q$ by $S_{c,p}x_q$. A linear term $c x_p$ may
be written directly as $C_c x_p$, and a constant $c$ as $C_c$.
Signs remain in $\sigma_{ji}$. All new rows are quadratic and use only
unit signed occurrences. Substituting the forced constants and scaled
wires restores each original source row exactly. This is a polynomial
identity, independently checked in the release, rather than a test on
sampled source zeros.

The binary normalization and the height-one construction compose. Every
source tuple has a uniquely specified extension to the new coordinates
at a zero. For the delivered normalization, its constant wires are
constant and its scaled wires are linear in the source variables; the
full inverse on zero sets therefore still has degree at most two.
In fact, if $L$ denotes this explicit extension, the stronger identity
\begin{equation}
 \mathcal H(L(x))=\sum_j R_j(x)^2
 \label{eq:forward-value}
\end{equation}
holds for every source tuple, not only for source zeros. The constant
and scaling rows vanish identically under $L$, each pair of row squares
in \eqref{eq:block-sos} contributes twice the original squared residual,
and all agreement and unit blocks vanish.

\subsection{Size accounting and the instantiated equation}

Let $L=\sum_j k_j$ be the number of signed occurrences after normalization.
The row blocks use $3L$ variables: two aliases and one $w$ per occurrence.
Let $d_2$ count old-variable labels appearing exactly twice among aliases,
and let $\delta_u\in\{0,1,2\}$ be the number of additional unit aliases
needed to make the unit group have size at least two. The exact count is
\begin{equation}
 N=3L+d_2+\delta_u.
 \label{eq:height-count}
\end{equation}
A block with $k$ occurrences contributes $k^2+2k$ distinct monomials:
$k$ quartic squares, $\binom{k}{2}$ quartic cross terms, $k$ cubic terms,
and $k+\binom{k}{2}$ quadratic terms. A $d$-vertex agreement cycle
contributes $2d$ monomials, and a $d$-unit anchor contributes $2d+2$.
These formulas are exact because the proof excluded collisions.

The instantiated source is the 519-integer, 321-row C222 system from
the incoming powered-saturation packet \cite{round14}, with its complete
quadratic residual list included unchanged. The artifact inventory and
measured counts appear in Section~\ref{sec:artifacts}. In particular,
the output contains every collected coefficient of $\mathcal H$; it
does not leave an input program, compiler choice, or coefficient oracle
to be supplied later.

The general construction changes the coefficient objective decisively:
height one is the least positive coefficient height of any nonzero
integer polynomial. This observation is only an optimum for that one
coordinate of the complexity problem. It makes no claim that the
resulting number of variables, support, or total description length is
minimal. The existing height-two construction \cite{heighttwo} has fewer
variables, and the existing native quartic has far fewer variables with
much larger coefficients. Those remain distinct tradeoffs.

\subsection{Logical consequence and precise novelty scope}

The construction is entirely finite. Its soundness and completeness use
identities of integer polynomials, finite dependency induction, and
nonnegativity of squares. These arguments are formalizable in ZFC.
Composing Theorem~\ref{thm:heightone} with the retained C222 calibration
therefore gives
\begin{equation}
 \ZFC\vdash
 \left[\forall z\in\mathbb Z^{\HeightOneVariables},\ \mathcal H(z)\ne0\right]
 \ \longleftrightarrow\ \Con(\ZFC).
 \label{eq:heightone-consistency}
\end{equation}
The height-one transformation does not introduce another unbounded
number-theoretic lemma. Its dependence on the inherited arithmetic
encoding is identified separately in Section~\ref{sec:calibration}.

More generally, quadratic circuit conversion followed by this theorem
reduces integer Diophantine solvability to multiquadratic quartics of
height one and constant coefficient one. This is a direct corollary of
the proved construction; it is not a claim to have newly established
the classical undecidability theorem. Searches of the frozen repository,
the current incoming packets, and selected literature did not locate
this particular block-and-cycle compiler. That bounded audit is not a
worldwide priority result. The concrete contribution is the proved
compiler, its independently checked fully instantiated consistency
polynomial, and its comparison with the supplied height-two endpoint.

\section{Eleven integers at degree 464}\label{sec:eleven}

This branch minimizes a different part of the complexity vector. It
retains eleven coordinates and the complete native source34 coefficient
tuple, but combines two incoming constructions in a way that substantially
reduces degree. The three-root coprime projection of \cite{coprime}
supplies the variable count. Powered saturation and supported-modulus
interpolation from \cite{round14} allow a much smaller scalar carrier.
Two elementary changes to the guard and the radical weights then reduce
the literal arithmetic cost.

\subsection{The complete equation}

The eleven independent integers are
\[
 j,C,z,h,r,v,w,T,c,n,s.
\]
The six fixed coefficients are denoted by
$\beta,K_0,M_C,M_F,A_0,B_0$. They are the unchanged source34 tuple,
with
\[
 \beta=2^{1064921}-1,\qquad A_0=2^{1113},\qquad B_0=2^{666104}.
\]
The complete bytes of $K_0,M_C,M_F$, as well as the other three
coefficients, are supplied with their bindings and finite reconstruction
recipe. In particular, $M_F$ already includes the source shift. None of
these six coefficients is a free parameter or an existential unknown.

Define the following polynomial abbreviations:
\begin{align}
 J&=6j+1,\qquad q=\beta J+1,\qquad W=12w+4,\nonumber\\
 X&=qW,\qquad Y=(2v+1)q^3,\qquad E=XY,\nonumber\\
 a&=E+Y,\qquad H=4a+3,\nonumber\\
 \Delta&=a^2+H=(a+2)^2-1,\label{eq:outer-one}\\
 F&=(K_0+W)C-z(q-1),\nonumber\\
 R&=(q(q-F)-C)(q^2-1)+(M_C+qM_F)J,\nonumber\\
 k&=R+1+hE,\qquad \eta=c-Yk,\label{eq:outer-two}\\
 D_{\rm err}&=(A_0F-q)^2+(B_0C-q)^2,\nonumber\\
 \mathcal B&=k^2(q^2-D_{\rm err})-4q^2\eta^2,\qquad D=3\mathcal B-4,\label{eq:budget}\\
 L&=XY^2k,\qquad D_t=1+L(L+k),\nonumber\\
 D_g&=1+D D_t s^2,\qquad P=D_tD_g.\label{eq:guard}
\end{align}
The new strong and carrier scales, and the auxiliary data, are
\begin{align}
 I&=2rc^2, & S&=srH,\label{eq:new-scales}\\
 D_f&=1+\Delta I^2, & Q&=\Delta^2 I^2, & \delta&=Q-1,\nonumber\\
 V&=D_f(c-R+4IT)-c,\nonumber\\
 C_y&=QV^2-1, & D_a&=\delta C_y,\label{eq:auxiliary}\\
 K&=X(3c-2X)+2, & Z_c&=HT+K,\nonumber\\
 A_c&=(Z_c^2+3H^2)(Z_c^3-19H^3).\label{eq:carrier}
\end{align}
For the centered equation set
\begin{equation}
 Z=S(n+QV)-A_c,\qquad
 u=4S^2D_f,\qquad b=S^2D_a,\qquad e=16S^2P.
 \label{eq:three-roots}
\end{equation}
Finally define
\begin{equation}
 \boxed{\displaystyle
 \mathcal E_{464}=
 \bigl[(Z^2-u-b-e)^2-4(ub+ue+be)\bigr]^2
       -64ubeZ^2=0.}
 \label{eq:E464}
\end{equation}
Equations \eqref{eq:outer-one}--\eqref{eq:E464} specify a single
ordinary integer polynomial. Abbreviations do not introduce coordinates,
and all exponents are fixed nonnegative integers. The literal circuit
contains 24 additions, 21 subtractions, and 77 multiplications: 122 live
binary operations in total. The direct companion replaces $Z$ by
$Sn-A_c$ and has exact degree 488 in 121 operations.

The universal identity underlying \eqref{eq:E464} is
\begin{equation}
 \mathcal E=
 \prod_{(\epsilon_1,\epsilon_2,\epsilon_3)\in\{-1,1\}^3}
 (Z-\epsilon_1\sqrt u-\epsilon_2\sqrt b-\epsilon_3\sqrt e).
 \label{eq:three-norm}
\end{equation}
It holds algebraically even for zero, negative, or dependent radicands.
One may prove it by successive quadratic conjugation, or work in the
free algebra generated by three formal square roots. The checker
verifies its entire polynomial expansion independently.

\subsection{The inherited source interface}

The old centered proof supplies four precise arithmetic interfaces,
whose statements and full ordinary source proofs are retained in the
package. They are used here in the order stated.

First, the disk inequality $\mathcal B\ge2$ yields the common strict
sign of $q,F,C,q-F-C$ and
\begin{equation}
 k\ne0,\quad D_{\rm err}<q^2,\quad
 \left|\frac{c}{k}-Y\right|<\frac12,\quad
 \sgn(c)=\sgn(kY).
 \label{eq:centered-strip}
\end{equation}
Second, the strong and auxiliary Pell equations, together with
$c^2\mid I$, recover the missing main Pell ordinate by absolute index
folding. Third, the \emph{pre-projection} part of the signed rank
argument yields
\begin{equation}
 p=R=2t+1,\quad t\ge24,\quad
 |Y|>\frac{(|X|-3)^t}{2}-1,\quad
 |H|>|X|(|X|-3)^t.
 \label{eq:rank-bound}
\end{equation}
It does not assume the main scalar projection. Fourth, after that
projection has been recovered, the source orientation and half-binomial
arguments regenerate a valid old native certificate and hence a finite
syntactic refutation. Using the fourth interface before proving the
projection would be circular; this construction does not do so.

For completeness, a canonical source witness from a refutation supplies
$a+2>1$ even, $R\equiv3\pmod4$, $c\equiv R\pmod4$, a positive
centered budget, and the old main projection. The fixed-pair interpolation
theorems described below apply to these canonical witnesses. These
retained source facts are substantial inherited mathematics, rather
than consequences of finite polynomial testing.

\subsection{Bootstrap without a hidden nonzero assumption}

The literal coefficient tuple gives, for every integer assignment,
\[
 q\equiv0\pmod2,\quad q\equiv2\pmod3,\quad
 |q|\ge\beta-1,\quad |X|\ge4|q|,\quad |Y|\ge|q|^3.
\]
In particular $H$ is odd and nonzero, and $\Delta>1$ with
$\Delta\equiv3\pmod4$. Writing $\ell=XY^2$ gives
$D_t=1+\ell(\ell+1)k^2\ge1$. It is odd, and it is one modulo three
because $\ell$ is zero or two modulo three.

\begin{lemma}[The reduced guard]\label{lem:guard}
The budget $\mathcal B$ is divisible by four for every integer
assignment. At a zero of \eqref{eq:E464}, one has
$S\ne0$, $s r c I\ne0$, $\mathcal B\ge2$, $D_g>0$, and $D_a>0$.
\end{lemma}
\begin{proof}
The integers $q,A_0,B_0$ are even. Therefore both $q^2$ and
$D_{\rm err}$ are divisible by four, as is the complete expression
for $\mathcal B$ in \eqref{eq:budget}. It follows that $D$ is divisible
by four and $D_g\equiv1\pmod4$. This makes the old external factor
four in $1+4DD_t s^2$ unnecessary.

If $S=0$, then $u=b=e=0$, and the polynomial is $A_c^8$.
Let $\pi(x)=(x^2+3)(x^3-19)$. Since
\[
 A_c=H^5\pi(T+K/H),
\]
it cannot vanish: the quadratic factor has no real root and the cubic
has no rational root. Thus $S\ne0$, so $s,r\ne0$.

At this point $c\ne0$ must not yet be assumed. If $c=0$, then
$I=Q=0$, $D_f=1$, $\delta=C_y=-1$, and $D_a=1$.
If $c\ne0$, then $|I|\ge2c^2$ and $D_f>|c|+1$. In the integer
expression $V=D_f(c-R+4IT)-c$, either the parenthesized integer is zero,
giving $|V|=|c|\ge1$, or $|V|\ge D_f-|c|>1$. Hence $Q>1$,
$C_y>0$, and $D_a>0$. Thus $D_a$ is positive in both cases.

If $D_g<0$, its radical in \eqref{eq:three-norm} is the sole nonzero
imaginary term: $D_f,D_a$ are positive, $D_t\ge1$, and $S\ne0$.
No conjugate factor can vanish. Since $D_g\equiv1\pmod4$, it follows
that $D_g>0$. If $\mathcal B\le1$, its divisibility by four gives
$\mathcal B\le0$, so $D\le-4$ and $D_g\le-3$, a contradiction.
Consequently $\mathcal B\ge2$.

The disk argument now proves \eqref{eq:centered-strip}. Because $Y$ is
a nonzero integer, that strip excludes $c=0$. Thus $I\ne0$ as well.
\end{proof}

\subsection{Separating three radicals and recovering the main root}

At a zero, divide the vanishing conjugate factor by $S$. After a choice
of signs, the rational number on its left is a signed sum
\begin{equation}
 2\sqrt{D_f}+\sqrt{D_a}+4\sqrt{P}.
 \label{eq:weights}
\end{equation}
All three radicands are positive odd integers. Indeed $I$ is even, so
$D_f\equiv D_a\equiv1\pmod4$, while $D_t,D_g$ are odd.
For each odd integer radicand, factor its square root as an odd integer
times the square root of its odd squarefree part. In any nontrivial
squareclass occurring in \eqref{eq:weights}, the coefficients have
distinct 2-adic valuations drawn from $1,0,2$. The term of least
valuation cannot cancel with the others. Linear independence of distinct
squarefree radicals then forces all three roots to be rational, hence
integral.

The identity
\begin{equation}
 \gcd(D_t,D_g)=\gcd(D_t,1+DD_t s^2)=1
 \label{eq:coprime}
\end{equation}
shows that the positive square $P=D_tD_g$ has two separately square
factors. Thus there exist integers $f,\tau,g,\widetilde y$ with
\[
 f^2=D_f,\quad \tau^2=D_t,\quad g^2=D_g,\quad
 \widetilde y^2=D_a.
\]
Put $y=\widetilde y/\delta$ initially in $\mathbb Q$.
It satisfies $QV^2-(Q-1)y^2=1$. The retained adjacent-coefficient
Pell descent makes $y$ an integer. Its hypotheses $Q>1$, integral
$V$, and square $Q=\Delta^2I^2$ have already been established.
The signed factor now also gives $S\mid A_c$.

For clarity, the main-root recovery occurs before scalar projection.
Set $B=|a+2|$, $C_*=|c|$, and $I_*=|I|$. Strong Pell classification
writes
\[
 |f|=X_B(m_0),\qquad I_*=Y_B(m_0),\qquad |f|>2I_*.
\]
Here $X_B(j),Y_B(j)$ are the usual integer coordinates of
$(B+\sqrt{B^2-1})^j$. Classification of the auxiliary norm and its
two residues gives an odd index $\nu$ and
$Y_B(\nu)\equiv\pm C_*\pmod{|f|}$. Fold that index modulo $2m_0$
and reflect to $0\le p\le m_0$. Since
$Y_B(p)+C_*\le2I_*<|f|$, no modular wrap is possible, and
$C_*=Y_B(p)$ with $p\ge1$.

The two residues used here are unchanged:
\[
 V\equiv-c\pmod f,\qquad V\equiv-R\pmod c.
\]
The second follows directly from
$V+R=\Delta I^2(c-R)+4ID_fT$ and $c^2\mid I$.
Lucas composition supplies the required rank divisibility with an even
multiplier $I/c^2$; no odd-multiplier normalization is assumed. The
retained pre-projection argument now applies and proves
\eqref{eq:rank-bound}.

\subsection{Powered divisibility replaces a large carrier}

The old construction used $S=IH^5$ to force $H\mid K$. Our smaller
$S=srH$ forces only $H\mid K^5$. The exact coefficient-size bounds
are strong enough to recover the same final projection from this weaker
statement.

Let
\[
 K_{\rm quad}=X(X+2ac)-1,\qquad \sigma=\sgn(c),\qquad
 \Pi=(\sigma X-2^R)(2^R\sigma X+1).
\]
The literal identity
\begin{equation}
 K=HXc-2K_{\rm quad}
 \label{eq:short-projector}
\end{equation}
is verified by expansion. Since $H\mid S\mid A_c$ and
$A_c\equiv K^5\pmod H$, we obtain $H\mid K^5$. Oddness of $H$
then gives $H\mid K_{\rm quad}^5$. The signed main Pell projection
identity supplied by the recovered main ordinate is
\[
 2^R K_{\rm quad}\equiv\Pi\pmod H.
\]
Raising this congruence to the fifth power gives $H\mid\Pi^5$.
Only powers of two, which are units modulo the odd $H$, are cancelled.

\begin{lemma}[The fifth power is smaller than the modulus]
Under the established pre-projection hypotheses, $|\Pi|^5<|H|$.
\end{lemma}
\begin{proof}
Put $x=|X|$. The unchanged actual coefficient tuple gives $x\ge2^{64}$,
far more than is needed. With $R=2t+1$ and $t\ge24$, elementary
estimates give
\[
 |\Pi|^5<2^{20t+20}x^{10},\qquad
 |H|>x(x-3)^t\ge\frac{x^{t+1}}{2^t}.
\]
The first inequality follows by bounding each of the two factors in
$\Pi$ by $2^{R+1}x$. The comparison needed for the two displayed
bounds is $x^{t-9}>2^{21t+20}$. It follows from
\[
 64(t-9)-(21t+20)=43t-596\ge436>0.
\]
\end{proof}

Divisibility and the strict height bound force $\Pi=0$. The second
factor of $\Pi$ is odd and therefore nonzero, so
$\sigma X=2^R$. The retained orientation and half-binomial proof can
now be used: it gives $X,c,k,Y>0$, the required strict native strip,
and a fresh valid old certificate. This completes soundness of the new
equation relative to the inherited native calibration.

\subsection{Completeness with a fixed strong pair}\label{sec:interpolation}

Begin with a canonical outer witness supplied by a finite syntactic
refutation in the retained source theorem. Its $\mathcal B$ is positive,
$D_t=\tau^2$, and $DD_t$ is a positive nonsquare: it is two modulo
three. Ordinary Pell existence supplies positive $g,s$ such that
\[
 g^2-DD_t s^2=1.
\]
This directly gives the reduced guard $D_g=g^2$. There is no need to
square this unit to create an extra factor four.

Choose one positive strong pair $I,f$ with $4sc^2\mid I$ and put
$r=I/(2c^2)$. Then $r>0$ and $2s\mid r$. Every prime divisor of
$M=sr$ divides $I$. Hold this strong pair, $r$, and $M$ fixed.

The supported-modulus interpolation theorem of \cite{round14} says
that every residue class of $T$ modulo any such $M$ is realized by an
integer auxiliary Pell witness $y$ for this fixed strong pair. It
applies to the same canonical even base and the same auxiliary chart
used here. Its proof does not depend on the choice of determinant or
centered-budget guard. The independent review checks the hypotheses
$a+2$ even, $R\equiv c\equiv3\pmod4$, and $4sc^2\mid I$ explicitly.

We recall the mechanism because the order of quantifiers matters.
Writing $I=Y_A(m_0)$ and $f=X_A(m_0)$, the exact period needed modulo
$f^2$ is $\lambda=4m_0f$. Set $d=\operatorname{lcm}(\lambda,4I)$ and
vary a starting auxiliary index through $\nu(h)=\nu_0+dh$.
The auxiliary coordinate
\begin{equation}
 T(h)=\frac{X_{\Delta I}(\nu(h))/(\Delta I)+c-f^2(c-R)}{4If^2}
 \label{eq:interp-coordinate}
\end{equation}
is an integer. Its denominator is part of the interpolation proof, not
part of the final polynomial. For every prime $p\mid I$, Pell addition
and valuation lifting give
\begin{equation}
 v_p\bigl(T(h_1)-T(h_2)\bigr)=v_p(h_1-h_2).
 \label{eq:isometry}
\end{equation}
The exact rank valuations give $v_p(d)=v_p(4I)$; the factor $f$ in
$\lambda$ cannot be omitted. For $p=2$, the first-coordinate difference
has valuation $2(v_2(\Delta I)+v_2(d(h_1-h_2)))-1$, strictly larger
than the ordinate contribution because the index difference is
divisible by four. Thus the two-adic case is included separately.
Equation \eqref{eq:isometry} induces a permutation of each supported
prime-power residue ring. Chinese remaindering of the corresponding
$h$ classes proves the full supported-modulus statement.

The polynomial $\pi(x)=(x^2+3)(x^3-19)$ has a root modulo every positive
integer. The inherited elementary proof treats the exceptional primes
$2,3,19$, uses lifting at the other primes, and combines the roots by
CRT. Since the canonical projection gives $H\mid K$, write
$\kappa=K/H$. Choose a residue with
\[
 \pi(T+\kappa)\equiv0\pmod M,
\]
and only then choose its interpolated auxiliary witness. It follows that
\[
 \frac{A_c}{S}=
 \frac{H^4\pi(T+\kappa)}{sr}\in\mathbb Z.
\]
Finally set
\begin{equation}
 n=\frac{A_c}{S}+2f+\delta y+4\tau g-QV.
 \label{eq:choose-n}
\end{equation}
One conjugate factor in \eqref{eq:three-norm} now vanishes, giving an
integer zero of $\mathcal E_{464}$.

The witness order is outer native data, guard Pell unit, one strong
pair, carrier residue, auxiliary $T,y$, and finally $n$. The strong
pair is never changed after the residue is prescribed. The new
construction preserves existence of zeros; it does not identify
old and new auxiliary tuples pointwise.

\subsection{Exact degree and paid arithmetic cost}

The principal degree ledger is
\begin{center}
\begin{tabular}{lrrrrrrrr}
\toprule
Expression &$I$&$S$&$D_f$&$Q$&$V$&$D_a$&$P$&$A_c$\\
Degree &3&8&18&30&23&106&94&35\\
\midrule
Expression &$u$&$b$&$e$&$Z_0=SQV$&$Z_1=Sn-A_c$\\
Degree &34&122&110&61&35\\
\bottomrule
\end{tabular}
\end{center}
The center uses the already computed $QV$. Its exact cancellation is
\begin{equation}
 b-Z_0^2=S^2(-QV^2-Q+1).
 \label{eq:center-cancel}
\end{equation}
Thus $\omega=(Z_0+Z_1)^2-b$ has degree at most 96. Rewriting the
identical polynomial as
\begin{equation}
 \bigl[(\omega-u-e)^2-4(ue+b(u+e))\bigr]^2
       -64ube(Z_0+Z_1)^2
 \label{eq:reduced-degree}
\end{equation}
shows that the unique highest-degree part inside the outer square is
$-4b_{122}e_{110}$, of degree 232. The last product has degree
$34+122+110+122=388$. Therefore the degree is exactly 464, with leading
form $16b_{122}^2e_{110}^2$.

The independently checked explicit coefficient
\begin{equation}
 [j^{260}w^{56}v^{64}C^4s^{12}r^{24}c^{32}h^{12}]\mathcal E_{464}
   =2^{484}3^{318}\beta^{260}>0
 \label{eq:leading464}
\end{equation}
proves attainment for the actual fixed tuple. The exponent sum is 464.
For the direct companion, $Z$ has degree 35 and the leading part is
$b_{122}^4$, of degree 488. An attained coefficient is
\[
 [j^{248}w^{64}v^{56}C^8s^8r^{40}c^{64}]\mathcal E_{488}
   =2^{480}3^{312}\beta^{248}>0.
\]
The degree certificate uses exact leading homogeneous arithmetic and
the cancellation identity, not a numerical evaluation of the polynomial
at a few points.

The delivered graphs count every live arithmetic node, authenticate all
six coefficient payloads, and reconstruct the coefficient recipe.
The 122-operation count is separate from the cancellation-aware degree
argument: an uncollected syntactic degree bound for the shortest graph
would overestimate the true degree.

\begin{theorem}[The eleven-integer equation]\label{thm:eleven}
For the six fixed coefficients supplied with the packet,
$\mathcal E_{464}$ has eleven unrestricted integer unknowns, exact
degree 464, and a 122-operation division-free evaluation circuit.
The retained source calibration and the arguments above prove
\[
 \ZFC\vdash
 [\forall x\in\mathbb Z^{11}\;\mathcal E_{464}(x)\ne0]
 \longleftrightarrow\Con(\ZFC).
\]
The same conclusion holds for the degree-488, 121-operation direct
companion.
\end{theorem}

\section{Calibration to syntactic consistency}\label{sec:calibration}

It is useful to separate the two sources because they certify the same
logical statement through different arithmetic representations.

\subsection{The direct C222 source for the height-one quartic}

The supplied \file{height_one/source/equation.json} fixes 519 integer
variables and 321 quadratic residuals $R_j$. It is copied byte for byte
from the centered direct endpoint in \cite{round14}. Its SHA-256 digest is
\begin{center}
\small\file{2a7dec9de97fdf80de1fdef4977b6cdf096bff5fc1dc7a586cdb5fede5333016}.
\end{center}
The inherited ordinary calibration is
\begin{equation}
 \ZFC\vdash
 \left[\exists x\in\mathbb Z^{519}\;\bigwedge_{j=1}^{321}R_j(x)=0\right]
 \longleftrightarrow
 \left[\text{there is a finite C222 refutation}\right]
 \longleftrightarrow\neg\Con(\ZFC).
 \label{eq:source-calibration}
\end{equation}
The source uses a finite numerical proof calculus, a direct arithmetic
encoding of its certificates, Pell and CRT arithmetic, integer sign
representations, exact monic eliminations, and an integral translation.
It does not use a simulation of K215 or a specialization of Jones's
universal polynomial.

The source logical endpoints are
\begin{quote}
\file{KEndpoint.c222_iff_zf_refutation}\\
\file{ConstructibleChoice.zf_refutation_iff_zfc_refutation}.
\end{quote}
These are inherited source interfaces. The direct arithmetic bridge
between the literal residuals and those endpoints is an ordinary proof,
not a theorem newly checked by Lean in this packet. The source's main
ordinary proof excerpt is retained in \file{height_one/source/section.tex};
additional arithmetic and calculus proof dependencies are retained in
\file{height_one/source/reading/}, with a reading order and exact source
hashes. This curated supplement is not a full ancestor compiler replay or
a transitive Lean dependency closure. In particular, the signed comparison
index theorem is an ordinary extension of the natural-coordinate Pell
criterion; the natural-coordinate theorem alone does not cover arbitrary
signed comparison coordinates.

The new coefficient compiler has an unconditional algebraic theorem
relative to this literal input system. It proves an explicit bijection
of zero sets and the exact value identity \eqref{eq:forward-value}.
Consequently it does not require another assumption about proof coding,
Pell completeness, or the consistency of any theory. Composing its
finite proof with \eqref{eq:source-calibration} yields
\eqref{eq:heightone-consistency}.

\subsection{The native source for the eleven-variable equation}

The eleven-variable construction inherits the coherent source34
coefficient tuple and the ordinary native/Pell/refutation calibration
retained by \cite{coprime}. The coefficient recipe fixes all six
integers by finite computations and payload hashes. It neither asks
whether ZFC is consistent nor searches for a refutation to choose a
coefficient. Its semantic route includes the earlier native encoding
and K215 calibration; it is distinct from the direct C222 route above.

The new proof establishes both reductions needed to reuse that
calibration. Starting from a zero, it proves nonzero conditions,
positivity, radical integrality, the main Pell ordinate, rank and ratio
bounds, the exact projection, and then reconstructs a valid old source
witness. Starting from a source witness, it chooses a new guard unit,
a fixed strong pair, an interpolated auxiliary residue, and the final
integer $n$. It is an equivalence of zero existence with regenerated
auxiliary witnesses. An equality of old and new polynomials is not
asserted.

\subsection{What is proved inside ZFC}

Each finite polynomial is an explicitly specified object of integer
arithmetic. ZFC proves the required facts about ordered integers,
finite sequences, elementary Pell theory, divisibility, and the finite
proof transformations used by the retained sources. Thus the combined
ordinary arguments prove the equivalences \eqref{eq:target} for the
stated polynomials. This formalizability claim is different from
providing an encoded derivation in a particular implementation of ZFC
or a transitive Lean/Rocq build.

No consistency hypothesis is needed to state or prove the reductions.
In the usual metamathematical reading, if ZFC is consistent then the
equations have no integer zeros. The report does not prove that premise
or the resulting nonvanishing statement outright. Equally, an integer
zero would, through the specified reductions, encode a finite syntactic
refutation. Numerical experiments cannot distinguish these alternatives.

\subsection{Priority and attribution}

The height-one block-and-cycle compiler was derived in this round and
reviewed independently within the AI-assisted workflow. Its elementary
energy identities are not claimed as new identities in mathematics.
The concrete coefficient-height improvement is measured against the
incoming height-two artifact, not against an exhaustive historical
classification of restricted-coefficient Diophantine equations.

The degree-464 equation combines the incoming coprime projection and
powered saturation/interpolation results, with the explicitly stated
guard and center refinements. These earlier mathematical contributions
remain credited. The provided independent reviews are separate agent
checks in this workflow; they are not external human endorsements or
formal-kernel certificates. Inherited notices and licenses are retained.

\section{Explicit artifacts and reproducibility}\label{sec:artifacts}

The archive is intended to be integrated as one research packet. The
new checker entry points use relative paths and can be run offline.
The complete polynomial data are supplied, rather than only instructions
to select a polynomial after a future computation.

\subsection{The height-one polynomial}

\file{height_one/output/polynomial.json} contains the ordered coordinate
metadata and every one of the 34,086 nonzero coefficients. A sparse term
$[c,[i_1,\ldots,i_k]]$ denotes
$c z_{i_1}\cdots z_{i_k}$, with zero-based coordinate indices and
repetition denoting powers. The empty index list denotes a constant.
The companion \file{polynomial.txt} prints the entire ordinary equation
in the same coordinate order.

\begin{center}
\begin{tabular}{lr@{\qquad}lr}
\toprule
Metric & Count & Degree of monomial & Count\\
\midrule
Original coordinates & 519 & 0 & 1\\
Original quadratic rows & 321 & 1 & 2\\
Binary constant wires & 96 & 2 & 20,853\\
Scaled-factor wires & 155 & 3 & 2,378\\
Normalized coordinates & 770 & 4 & 10,852\\
Normalized rows & 572 & Total & 34,086\\
\midrule
Signed term occurrences & 2,378 & Coefficient $+1$ & 21,934\\
Equality-cycle dummies & 297 & Coefficient $-1$ & 12,152\\
Final coordinates & 7,431 & Unit aliases & 1,208\\
\bottomrule
\end{tabular}
\end{center}
All 519 input coordinates, all 770 normalized coordinates, and all 7,431
output coordinates occur. The exact dimension identity is
$3\cdot2378+297=7431$. There are 52 singleton alias groups; they
require no agreement cycle. The unit group already has more than two
members and needs no padding dummy.

The complete correspondence certificate records every constant and
scaling definition, signed occurrence, alias pair, product witness,
agreement cycle, and unit-path index. It is sufficient to reconstruct
the lifting and projection maps without importing the original proof
compiler. The independent verifier imports neither the new producer nor
upstream arithmetic code. It checks all 321 source-row pullbacks as
exact polynomial identities and independently expands the full
$2\mathcal H=\sum Q_i^2$ certificate.

The release also checks total and individual degrees, the coefficient
alphabet, constant term, variable occurrence, and exact equality of the
JSON and plaintext polynomials. Two actual modified-artifact checks
flip a coefficient and remove the unit-anchor constant; both are
rejected by the full verifier at the expanded-polynomial comparison.
Five deterministic signed evaluations supplement the exact checks;
they are not the proof of a universal identity.

\subsection{The eleven-variable circuits}

The two circuit JSON files under \file{eleven/} list integer leaves,
eleven coordinate leaves, six authenticated fixed-coefficient leaves,
and binary arithmetic nodes. Their output register fixes the complete
polynomial. \file{eleven/coefficients/} contains all six compressed
integer payloads, byte-order and binding information, and the finite
source recipe. A coefficient leaf therefore denotes a literal fixed
integer, not an unbound parameter.

The independent checker verifies the universal three-radical norm
identity, the actual graph definitions, dependency order, liveness,
variable set, arithmetic count, exact degree cancellations, and attained
leading coefficients. It authenticates the whole coefficient tuple and
reconstructs the native recipe. The proof review separately checks the
unbounded logical implications, including the case $c=0$, the reduced
guard, the exact Pell period, the supported-modulus hypothesis, and
the order of witness selection.

\subsection{Running the checks}

From the extracted packet root, the supported release command is
\begin{quote}
\file{python3 verify_all.py}
\end{quote}
It runs the two independent verifiers in a temporary copy, compares
deterministic regeneration with the shipped data, and checks the release
manifest. Both new compilers and their exact verifiers use only the
Python standard library; no SymPy, theorem prover, or network access
is required. The release runs ordinary and optimized Python checks.
The new verifiers use explicit exceptions and authenticate the rebuilt
coefficient values in both modes, even when assertions in an inherited
recipe are disabled.

To regenerate the PDF, run two passes, or use \file{latexmk}, from
\file{report/}. The LaTeX source and all included sections are supplied.
PDF bytes may differ with TeX versions, metadata, or timestamps;
polynomial JSON regeneration is byte-for-byte deterministic.

The included receipts record actual successful executions. They certify
the finite artifact checks described above. They do not constitute a
satisfiability search, a proof of Con(ZFC), or an automatically generated
derivation of the inherited unbounded arithmetic calibration.

\section{Further research questions}\label{sec:further}

\subsection{Reduce dimension while retaining height one}

The general compiler deliberately assigns two fresh aliases and one
product witness to every signed occurrence. This gives a transparent
collision proof but is not a minimal allocation. Several precise
optimization problems follow. Which linear occurrences can remain
unpadded while preserving disjointness between the quadratic agreement
terms and the row-block terms? Can two occurrences share a product
witness without introducing another zero or a coefficient of magnitude
two? Which agreement classes can be represented by a smaller positive
semidefinite form with the same integer kernel and coefficient alphabet?
Any optimization should supply both its zero-set proof and a complete
collected-coefficient audit.

The current binary normalization uses 96 constant wires and 155 shared
scaled-factor wires. Alternative addition chains, signed-digit constant
representations, and selective use of a different 519- or 517-variable
source may reduce the normalized occurrence count before the height-one
compiler is applied. The correct objective is the final count
$3L+d_2+\delta_u$, not the smallest intermediate source dimension alone.

\subsection{Compile the low-variable circuit directly}

Another route is to quadratize the 122-operation circuit and then apply
the height-one compiler. Its native coefficients have very large bit
lengths, so explicit coefficient construction may dominate the resulting
variable count. A meaningful comparison must include those constant
wires. Treating the six source ports as free constants in a circuit
count and then omitting their cost in a height-one encoding would not
be a fair comparison.

\subsection{Lower the eleven-variable degree further}

The degree-464 leading form is $16b_{122}^2e_{110}^2$. Thus the fused
guard radical and the auxiliary radical are the present bottleneck.
Reducing only an expression that lies below those two leading degrees
will not lower the final degree. Possible targets include a smaller
budget guard preserving the unconditional parity information, a
different coprime fusion, or a carrier that continues to exclude both
$r=0$ and $s=0$ with less degree.

The supported-modulus theorem must continue to be applied with a fixed
strong pair. A proposed carrier change is incomplete if it chooses a
residue first and then silently changes that pair to realize the residue.
Likewise, a weakened projection must compare the entire powered
divisibility product with the actual modulus. The generous authentic
coefficient bound may permit other exponents or smaller scales, but
each such claim needs a new quantitative proof.

\subsection{Formalize the local new transformations first}

The height-one compiler offers a relatively compact formalization
target: finite sparse polynomial identities, a graph-cycle kernel,
nonnegativity, and an explicit zero-set bijection. A proof assistant
could first verify the generic compiler and then check the literal
occurrence certificate. This would avoid formalizing C222 arithmetic
again merely to validate the coefficient transformation.

For the eleven-variable equation, useful separable targets are the
guard divisibility lemma, three-root separation with weights $2,1,4$,
the powered projection bound, and the supported-modulus interpolation
interface. The latter needs explicit attention to the exact period
$4m_0f$ and the two-adic valuation step. An eventual end-to-end result
must connect these local proofs to the concrete source equation and
its internal consistency statement, preserving the exact coefficient
tuple throughout.

\subsection{Define a reproducible simplicity objective}

The two new endpoints make the tradeoff visible. Height one is attained
with thousands of variables; eleven variables are attained with large
fixed coefficients and higher degree. A next comparison could minimize
a specified encoded byte length, weighted arithmetic cost, or
lexicographic tuple of degree, dimension, height, and support. Without
such a choice, a global claim of ``the simplest equation'' is not a
well-defined conclusion of this packet. What is established here is
two concrete, fully specified improvements and the proofs needed to
evaluate them.


\begin{thebibliography}{9}
\bibitem{repo}
Vladimir Reshetnikov and ProveIt contributors,
\emph{ProveIt: Diophantine consistency research}, snapshot
\file{45b399ec51dfbecb891dd51d8a9e4b6b7c361ac6}, 24 October 2026.
\href{https://github.com/VladimirReshetnikov/ProveIt/tree/45b399ec51dfbecb891dd51d8a9e4b6b7c361ac6/Computability/DiophantineConsistency}{Immutable repository tree}.

\bibitem{round14}
ProveIt incoming research packet,
\emph{Diophantine consistency equations with powered saturation and
factorial carriers, Round 14}, 24 October 2026.
Repository archive
\file{ConZFC_Powered_Saturation_Factorial_Carriers_Round14_2026-10-24.zip};
the source tuple, direct 519-variable system, and selected ordinary proof
dependencies are identified in the supplied provenance files.

\bibitem{coprime}
ProveIt incoming research packet,
\emph{Diophantine consistency equations by coprime projection},
24 October 2026. Repository archive
\file{ZFC_Diophantine_Coprime_Projection_2026-10-24.zip}, particularly
\file{integer/PROOF.md} and its retained centered source proofs.

\bibitem{heighttwo}
ProveIt incoming research packet,
\emph{Signed powers and coefficient height two}, 24 October 2026.
Repository archive \file{ConZFC_Signed_Powers_Height_Two_2026-10-24.zip},
\file{height_two/PROOF.md} and the \file{anchored/} endpoint.

\bibitem{bayer}
J. Bayer, M. David, M. Ha\ss ler, Y. Matiyasevich, and D. Schleicher,
\emph{Diophantine equations over $\mathbb Z$: universal bounds and
parallel formalization}, arXiv:2506.20909v3, 22 October 2026.
\url{https://arxiv.org/abs/2506.20909}.

\bibitem{sun}
Z.-W. Sun,
\emph{Ten unknowns for Hilbert's tenth problem over the integers},
arXiv:2609.01594v2, 2 September 2026.
\url{https://arxiv.org/abs/2609.01594}.
\end{thebibliography}
\end{document}
